% !TEX root = ../main.tex
\section{Conclusions and future work}
In this work, we provide a large-scale characterization of hosting and security practices of public-facing \gls{LDAP} servers. 
Our findings show a diverse \gls{LDAP} ecosystem with 15.3k hosting organizations, including many self-managed solutions. This diversity also increases the complexity of maintenance. We find evidence for outdated software, deprecated \gls{TLS} configurations, and vulnerable certificates, which supports the idea that network operators and \gls{LDAP} administrators face difficulties in keeping their systems secure.
To address these issues, we recommend that network operators adopt systematic monitoring and upgrade strategies, isolate deprecated instances, and follow established best practices for \gls{TLS} and certificate management.
%We make our code publicly available under a permissive license.
%We presented an analysis of public-facing \gls{LDAP} servers. We developed a methodology to study the networks where \gls{LDAP} servers are hosted.
%We collected and formed a dataset that we use to recognize \gls{LDAP} versions running on different OSes and determine their adopted security practices.
%We found that the hosting ecosystem of \gls{LDAP} servers is vast and distributed across 15.3k different organizations in 9.4k different ASes, without a dominant hoster as previously seen on the Web.
%We report on deprecated software running on large hosters such as Comcast, Amazon, and Hetzner. 
%We also report on servers that use certificates with compromised keys.
%Furthermore, our dataset, measurement, and analysis code are publicly available.
We aim to raise awareness among researchers and network operators about a more secure ecosystem, and we invite the community to extend this research by conducting qualitative or quantitative surveys on server maintenance.
